\documentclass[10pt,a4paper]{amsart}
\usepackage[margin=19mm]{geometry}
\usepackage{amsmath,amssymb,mathtools}
\usepackage[scr=rsfs,cal=euler]{mathalfa}
\usepackage{xfrac}
\usepackage[unicode,hidelinks,bookmarks=false]{hyperref}
\newcommand{\ie}{i.e.\ }
\newcommand{\cf}{cf.\ }
\newcommand{\resp}{respectively\ }
\newcommand{\Subsection}{\subsection}
\providecommand{\nobreakdash}{\nobreak\hbox{-}}
\setlength{\emergencystretch}{2em}
\multlinegap0pt
\usepackage{enumerate}
\usepackage{amssymb}
\usepackage{esint}
\usepackage{xcolor}
\usepackage{mathtools}
\usepackage{bm}
\usepackage{tikz}
\usepackage{float}
\newtheorem{theorem}{Theorem}
\newcommand\bR{\mathbb{R}}
\newcommand\cC{\mathcal{C}}
\newcommand\cF{\mathcal{F}}
\newcommand\cL{\mathcal{L}}
\newcommand\ep{\varepsilon}
\title{A Multiplicative Neural Network Architecture:\\ Locality and Regularity of Approximation}
\author{Hee-Sun Choi, Beom-Seok Han}
\date{Source: arXiv:2602.06374v2 (CC BY 4.0)}
\begin{document}
\maketitle

\noindent\emph{Source and licence.} A Multiplicative Neural Network Architecture:\\ Locality and Regularity of Approximation. Version arXiv:2602.06374v2 (CC BY 4.0), distributed by arXiv under the Creative Commons Attribution 4.0 International licence, http://creativecommons.org/licenses/by/4.0/, as recorded in the arXiv licence metadata for this exact version and shown on the abstract page https://arxiv.org/abs/2602.06374v2. Full licence notice: \texttt{licenses/arxiv-2602.06374-CC-BY-4.0.txt}.\par\smallskip

\noindent\textit{Editorial scope.} Counted unit: the article's theorem theorem (source0.tex lines 360--381). The article's complete original proof of it is delivered with it (source0.tex lines 588--737, one \texttt{proof} environment of 150 lines, which the article places away from the statement). No mathematics is written, repaired or re-derived: the statement and the proof are the article's own lines, in the article's own notation, and no retained line is substituted or re-flowed.  No further source-local context is needed: every cross-reference of every retained line resolves inside the two delivered spans.  Nothing was omitted from any retained span, so the package carries no omission record.  The retained lines cite 1 works, whose entries are transcribed verbatim from the article's own bibliography source in the e-print and delivered with short editorial labels recorded in the item's reference mapping.  No retained line contains a figure environment, an image-inclusion command or drawing code, so no image is delivered and the package declares no asset dependency.  Editorial formatting only, in addition: an AMS \texttt{amsart} preamble that restates the article's own declarations and macros used by the retained lines, namely the environments \texttt{theorem} on separate counters, together with the standard packages those lines need.  The retained environments print in this document as Theorem 1 in order of appearance, whereas the article prints the same items with the numbering of its own full document.  The complete native source of the article as distributed in the arXiv e-print for this exact version is bundled unchanged as \texttt{sources/194095/source0.tex}.
\begin{theorem}
\label{main theorem}
Let $\gamma\geq0$ and $p\in (1,\infty)$. Assume that $\sigma:\bR \to \bR$ is a nontrivial function such that $\sigma\in H_{p}^{\gamma}(\bR)$. Suppose $S$ is a set of functions $F:\bR^{m}\to \bR$ of the form
\begin{equation}
\label{eq:approximation function}
F(x) = \sum_{j = 1}^{n}\alpha_{j} \prod_{i = 1}^{m}\sigma\left( w^{T}_{ij} x + b_{ij} \right),
\end{equation}
where 
$n\in \{1,2,\dots,\}$, $\alpha_{j}\in\bR$,  $w_{ij}$, $x\in \bR^{m}$, $b_{ij}\in \bR$, and $
\begin{bmatrix}
w_{1j}~w_{2j}~\cdots~w_{mj}
\end{bmatrix}^{T}$
is an invertible matrix. Then, $S$ is dense in $H_p^{\gamma}(\bR^{m})$. In other words, given any $f\in H^{\gamma}_p(\bR^{m})$ and $\ep>0$, there exists the function $F$ of the form \eqref{eq:approximation function} such that
\begin{equation*}
\| F - f \|_{H_p^{\gamma}(\bR^{m})} \leq \ep.
\end{equation*}
In particular, if $\gamma - \frac{m}{p} >0 $, then $f\in \cC^{\gamma-\frac{m}{p}}(\bR^{m})$ and
\begin{equation*}
\| F - f \|_{\cC^{\gamma-\frac{m}{p}}((\bR^{m}))} \leq \ep,
\end{equation*}
where $\cC^{\gamma-\frac{m}{p}}(\bR^{m})$ is the Zygmund space.
\end{theorem}

\begin{proof}[Proof of Theorem \ref{main theorem}]
% Since $H_{p}^{\gamma}(\bR^{m})\cap L_{1}(\bR^{m})$ is dense in $H_{p}^{\gamma}(\bR^{m})$, we may assume that $F\in H_{p}^{\gamma}(\bR^{m})\cap L_{1}(\bR^{m})$. 
First we consider the case $\sigma\in C_{c}^{\infty}(\bR)$ is nontrivial. Let $S$ be a set of functions of the form \eqref{eq:approximation function}. 

First we show that $S$ is a linear subspace of $H_p^\gamma(\bR^{m})$. If we choose $f$ and $g$ in $S$, then for any $a,b\in\bR$, we have
\begin{equation*}
\begin{aligned}
af(x) + bg(x)
&= a\sum_{j = 1}^{n_{1}}\alpha_{j} \prod_{i = 1}^{m}\sigma\left( w^{T}_{ij} x + b_{ij} \right) + b\sum_{j = 1}^{n_{2}}\bar\alpha_{j} \prod_{i = 1}^{m}\sigma\left( \bar w^{T}_{ij} x + \bar b_{ij} \right) \\
&= \sum_{j = 1}^{n_{1} + n_{2}}c_{j} \prod_{i = 1}^{m}\sigma\left( \tilde w^{T}_{ij} x + \tilde b_{ij} \right),
\end{aligned}
\end{equation*}
where $c_{j} = a\alpha_{j}1_{1\leq j \leq n_{1}} +b\bar \alpha_{j-n_{1}}1_{n_{1} < j \leq n_{2}}$, $\tilde w_{ij} = w_{ij}1_{1\leq j \leq n_{1}}+\bar w_{i(j-n_{1})}1_{n_{1} < j \leq n_{2}}$, and $\tilde b_{ij} = b_{ij}1_{1\leq j \leq n_{1}}+\bar b_{i(j-n_{1})}1_{n_{1} < j \leq n_{2}}$

Thus, $S$ is a linear space. Next, observe that $S\subset H_p^\gamma(\bR^{m})$. Indeed,
\begin{equation*}
\begin{aligned}
\left\|   \prod_{k = 1}^{m} \sigma(w_{k}^{T}\cdot + b_{k})   \right\|_{H_{p}^{\gamma}(\bR^{m})}
& \leq N\left\|   \prod_{k = 1}^{m} \sigma(M_{k})   \right\|_{H_{p}^{\gamma}(\bR^{m})} 
\leq N\left\|    \sigma   \right\|_{H_{p}^{\gamma}(\bR)}^{m}<\infty,
\end{aligned}
\end{equation*}
where $w_{k}\in \bR^{m}$, $b_{k}\in \bR$, and 
\begin{equation*}
M_{k}(x) = x^{k}
\end{equation*} 
Note that the first inequality follows from the fact that the linear transformations only changes the norm constant times. In the case of second inequality, notice that
\begin{equation*}
\begin{aligned}
  (1-\Delta)^{\gamma/2}\prod_{k = 1}^{m} \sigma(x^{k})  
&= \int_{\bR^{m}} e^{2\pi i \xi \cdot x}(1+|\xi|^{2})^{\gamma/2}\int_{\bR^{m}} e^{-2\pi i \xi \cdot y} \prod_{k = 1}^{m}\sigma(y^{k}) dy d\xi  \\
&= \int_{\bR^{m}} e^{2\pi i \xi \cdot x} \mu(\xi^{1},\xi^{2},\cdot,\xi^{m}) \hat g(\xi) d\xi \\
&= \cL g(x)
\end{aligned}
\end{equation*}
where 
\begin{equation*}
\mu(\xi^{1},\xi^{2},\cdot,\xi^{m}) := \frac{(1+|\xi|^{2})^{\gamma/2}}{\prod_{k = 1}^{m} (1+|\xi^{k}|^{2})^{\gamma/2}}, \quad g(x) = \prod_{k = 1}^{m}\left[ \left( \left( 1-\frac{d^{2}}{dx^{2}} \right)^{\gamma/2}\sigma \right)(x^{k}) \right]
\end{equation*}
and 
\begin{equation*}
\cL(\cdot) = \cF^{-1}(m \cF(\cdot)).
\end{equation*}
is a zeroth-order pseudo-differential operator. Thus, 
\begin{equation*}
\begin{aligned}
\left\|   \prod_{k = 1}^{m} \sigma(M_{k})   \right\|_{H_{p}^{\gamma}(\bR^{m})} 
&= \left\|   (1-\Delta)^{\gamma/2} \prod_{k = 1}^{m} \sigma(M_{k})   \right\|_{L_{p}(\bR^{m})} \\
&= \left\| \cL g   \right\|_{L_{p}(\bR^{m})} \\
&\leq N\left\|  g   \right\|_{L_{p}(\bR^{m})} \\
&= N\left\| \prod_{k = 1}^{m}\left[ \left( (1-D^{2})^{\gamma/2}\sigma \right)(M_{k}) \right]   \right\|_{L_{p}(\bR^{m})} \\
&= N\left\| \left[ \left( (1-D^{2})^{\gamma/2}\sigma \right) \right]   \right\|_{L_{p}(\bR)}^{m} \\
&= N\left\| \sigma   \right\|_{H_{p}^{\gamma}(\bR)}^{m}.
\end{aligned}
\end{equation*}



Now, we prove that the space $S$ is dense in $H_p^\gamma(\bR^{m})$. To apply a proof by contradiction, assume that there is $f_0\in H_p^\gamma(\bR^{m})\setminus \bar S$. Then, by the Hahn-Banach theorem, there is a bounded linear functional $\Lambda$ on $H_p^\gamma(\bR^{m})$ such that $\Lambda f = 0$ for $f\in \bar S$ and 
\begin{equation}
\label{contradiction point}
\Lambda f_0\neq 0.
\end{equation}
% (by Theorem 5.19 of RCA.)
Then, by the Riesz representation theorem, there exists nonzero $h\in H_q^{-\gamma}(\bR^{m})$ with $q = p/(p-1)$ such that for any $f\in H_p^\gamma(\bR^{m})$,
\begin{equation*}
\Lambda f = (f,h).
\end{equation*}
For example, $(f,h) = \int_{\bR^{m}} f(x)h(x) dx$ if $\gamma = 0$.


For $\ep>0$ and $y = (y^{1},y^{2},\dots,y^{m})\in\bR^{m}$, set
\begin{equation}
\label{mollifier}
\xi (x) := \frac{1}{\| \sigma \|_{L_1(\bR)}^{m}}\prod_{i = 1}^{m}\sigma(x^{i})\quad\text{and}\quad 
\xi_{\ep}(x) := \ep^{-m}\xi\left( x/\ep \right).
\end{equation}
Since $\sigma\in C_c^\infty(\bR)$ is nontrivial, $\| \sigma \|_{L_{1}(\bR)}\neq 0$, and thus $\xi$ is well defined. 
Observe that the function
$x\mapsto \xi_\varepsilon(y-x)$ belongs to $S$ for each $y\in\mathbb{R}^m$. Indeed, for each $y\in\mathbb{R}^m$,
\[
\xi_\varepsilon(y-x)
=\varepsilon^{-m}\prod_{i=1}^m \sigma\!\left(\frac{y_i-x_i}{\varepsilon}\right)
\]
is of the form \eqref{eq:approximation function} with $n=1$, where 
\[
\begin{bmatrix}
w_{1j}~w_{2j}~\cdots~w_{mj}
\end{bmatrix}^{T} = W =\mathrm{diag}\!\left(\frac1\varepsilon,\dots,\frac1\varepsilon\right),
\]
and $b_i=y_i/\varepsilon$. Since $W$ is invertible, $\xi_\varepsilon(y-\cdot)\in S$ for all $y\in\mathbb{R}^m$.


Define
\begin{equation*}
h_{\ep}(y) 
:= \left(\xi_{\ep}(y - \cdot),h\right).
\end{equation*}
Then, $h_\varepsilon(y)
=(\xi_\varepsilon(y-\cdot),h)=0$ for all $y\in\mathbb{R}^m$,
and hence $h_\varepsilon\equiv0$.
On the other hand, since 
\begin{equation}
\label{mollification}
\|h-h_\varepsilon\|_{H_q^{-\gamma}(\mathbb{R}^m)} \to 0
\end{equation}
(e.g., see \cite[Theorem 13.9.2]{krylov2024lectures}) and $h_\varepsilon\equiv0$ for all $\varepsilon>0$, it follows that
\[
\|h\|_{H_q^{-\gamma}(\mathbb{R}^m)}=\|h-h_\varepsilon\|_{H_q^{-\gamma}(\mathbb{R}^m)}\to0,
\]
and hence $h=0$.
This contradicts \eqref{contradiction point} and therefore $S$ is dense in
$H_p^\gamma(\mathbb{R}^m)$.


Now we consider the case $\sigma\in H_{p}^{\gamma}(\bR)$. Let $\ep>0$. Since $C_{c}^{\infty}(\bR)$ is dense in $H_{p}^{\gamma}(\bR)$, there exists $\{\sigma_{\ell}\}\subset C_{c}^{\infty}(\bR)$ such that 
\begin{equation}
\label{sigma approximation}
\|\sigma_{\ell} - \sigma\|_{H_{p}^{\gamma}(\bR)}\leq 1/\ell
\end{equation}
By the previous step, there exists $F_{\ell}$ such that
\begin{equation*}
F_{\ell} = \sum_{j = 1}^{n}\alpha_{j} \prod_{i = 1}^{m}\sigma_{\ell}\left( w^{T}_{ij} x + b_{ij} \right)
\end{equation*}
and 
\begin{equation*}
\| F_{\ell} - f \|_{H_{p}^{\gamma}(\bR^{m})} \leq \ep.
\end{equation*}
Note that
\begin{equation*}
\left\|F - F_{\ell} \right\|_{H_{p}^{\gamma}(\bR^{m})} \leq N/\ell^m.
\end{equation*}
Indeed,
\begin{equation*}
\begin{aligned}
\left\|F - F_{\ell} \right\|_{H_{p}^{\gamma}(\bR^{m})}
&=\left\| \sum_{j = 1}^{n}\alpha_{j} \left(\prod_{i = 1}^{m}\sigma\left( w^{T}_{ij} \cdot + b_{ij} \right) - \prod_{i = 1}^{m}\sigma_{\ell}\left( w^{T}_{ij} \cdot + b_{ij} \right) \right)\right\|_{H_{p}^{\gamma}(\bR^{m})} \\
&\leq N\| \sigma - \sigma_{\ell} \|_{H_{p}^{\gamma}(\bR)}^{m} \\
% &\leq NM^{m}\ep/2 \\
&\leq N/\ell^m.
\end{aligned}
\end{equation*}
Thus,
\begin{equation*}
\| F - f \|_{H_{p}^{\gamma}(\bR^{m})} 
\leq \|F - F_{\ell}\|_{H_{p}^{\gamma}(\bR^{m})}
+ \|F_{\ell} - f\|_{H_{p}^{\gamma}(\bR^{m})} \leq N/\ell^m+\ep.
\end{equation*}
Since $\ell=1,2,\dots$ is arbitrary, the theorem is proved.
\end{proof}

\begin{thebibliography}{99}
\bibitem[krylov]{krylov2024lectures} N.~V. Krylov. \newblock {\em Lectures on elliptic and parabolic equations in Sobolev spaces}, volume~96. \newblock American Mathematical Society, 2024.
\end{thebibliography}
\end{document}
